\subsection{Parámetros revisados}

\begin{itemize}
  \item Desde el 12 de noviembre todo se elige con validación cruzada interna dentro del entrenamiento y la validación se mide una sola vez; las cifras del protocolo anterior no se usan.
  \item Las cifras finales se recalcularon con las 76 vacas del bucket actualizado el 5 de noviembre.
  \item Se probó foto original, foto sin fondo (segmentada con Mask R-CNN) y recorte de anca y cola; la foto sin fondo se adoptó como entrada estándar.
  \item De una regresión ridge fija a 25 recetas con selección anidada; agregar más candidatos empeoró el resultado.
  \item Estiramiento de la salida y pesos por rango, para corregir que las predicciones se concentren en el centro.
  \item El ajuste fino con varias semillas promediadas se lanzó, pero no se terminó por su costo de cómputo (unas 18 horas).
  \item Pesos públicos entrenados con fotos de anca de vacas lecheras (Modelo 4).
  \item De regresión a clasificación por cuartos de punto con valor esperado (Modelo 5).
  \item En la salida redondeo a 0.25, aplicado igual a todos los modelos y reportado siempre junto a la salida continua.
  \item El Modelo 5 se volvió a entrenar desde cero el 18 de noviembre, y la lectura de su salida (valor esperado redondeado, valor esperado continuo o clase más probable) se elige dentro de la validación de cada pliegue. Sus cifras en este documento son las de esa corrida.
\end{itemize}
